\chapter{Imatges i factoritzacions}\label{cap:imatges}

\begin{objectius}
\apartat{Prerequisits} Del capítol~\ref{cap:relacions}: parell nucli, quocient, epimorfisme regular i l'axioma d'efectivitat. Del capítol~\ref{cap:parts}: bimorfismes, imatge inversa, el graf d'una aplicació i l'equivalència entre injectivitat i monomorfisme. Del capítol~\ref{cap:valors}: l'exhaustivitat dels epimorfismes sobre punts. L'axioma d'elecció no s'usa en aquest capítol.
\apartat{Objectius} En acabar el capítol, el lector hauria de poder:
\begin{itemize}
  \item definir la imatge d'una aplicació com a part mínima per la qual es factoritza, construir-la com a quocient del parell nucli i demostrar-ne la unicitat;
  \item demostrar, sense elecció, que el factor monomòrfic de la factorització és un monomorfisme, i que tot epimorfisme és regular;
  \item demostrar que els epimorfismes són estables per canvi de base i explicar en quin sentit la teoria és regular;
  \item caracteritzar monomorfismes, epimorfismes i isomorfismes per la degeneració d'un dels dos factors de la factorització imatge;
  \item definir la imatge directa d'una part i demostrar l'adjunció amb la imatge inversa;
  \item compondre relacions i llegir la composició per elements generalitzats, amb el paper dels recobriments epimòrfics;
  \item identificar les aplicacions amb les relacions funcionals i totals, i la composició de relacions amb la d'aplicacions.
\end{itemize}
\end{objectius}

El capítol precedent va lligar les relacions d'equivalència amb els
epimorfismes regulars, i va mostrar que tot epimorfisme regular és el
quocient del seu parell nucli. Aquestes eines tenen una conseqüència
immediata que encara no hem explotat: tota aplicació, sigui quina sigui,
es descompon de manera essencialment única en una part exhaustiva seguida
d'una part injectiva. La part injectiva és la \emph{imatge} de
l'aplicació.

Aquest és un capítol de consolidació. No hi afegim cap axioma: tot el que
segueix es dedueix dels productes, els pullbacks i els coequalitzadors ja
disponibles, i de l'axioma d'efectivitat del
capítol~\ref{cap:relacions}. Tampoc no hi fem servir l'axioma d'elecció:
en el punt on seria més temptador ---la prova que el factor monomòrfic de
la factorització imatge és un monomorfisme--- n'hi ha prou amb la
separació global i amb l'exhaustivitat dels epimorfismes sobre punts
(observació~\ref{obs:factoritzacio-eleccio}). Una conseqüència
d'aquesta via és que tot epimorfisme resulta regular sense elecció, cosa
que al capítol anterior només havíem obtingut amb ella. El resultat és
una manera uniforme de parlar
de la imatge d'una aplicació i de la imatge directa d'una part, que
completa la imatge inversa introduïda al capítol~\ref{cap:parts} i precisa
en quin sentit la teoria obtinguda és regular. Com sempre, definirem les
nocions sobre elements generalitzats i en recuperarem la lectura clàssica
com a cas particular.

\section{La imatge d'una aplicació}\label{sec:imatge}

Sigui \(f\colon A\to B\) una aplicació. El seu parell nucli
\(\operatorname{Nuc}(f)\) (definició~\ref{def:kernel-pair}) és una relació
d'equivalència sobre \(A\) (proposició~\ref{prop:kernel-equivalencia}):
identifica els elements generalitzats de \(A\) que \(f\) envia al mateix
lloc. És natural, doncs, prendre'n el quocient i comparar-lo amb \(B\).

\begin{definicio}[Imatge d'una aplicació]\index{Imatge d'una aplicació}\label{def:imatge}
Sigui \(f\colon A\to B\). Una \emph{imatge} de \(f\) és una part
\(m\colon I\hookrightarrow B\) per la qual \(f\) es factoritza com
\[
f=m\circ e,
\qquad
e\colon A\twoheadrightarrow I,
\]
amb \(e\) un epimorfisme, i que és \emph{mínima} amb aquesta propietat: si
\(m'\colon I'\hookrightarrow B\) és una part i \(f\) es factoritza a
través de \(m'\), aleshores \(m\) es factoritza a través de \(m'\).
\end{definicio}

\begin{notacio}[La imatge]\label{not:imatge}
Escriurem \(\operatorname{Im}(f)\) per a la part \(m\colon
I\hookrightarrow B\) de la definició anterior, i
\(\operatorname{Imatge}(m,e,f)\) per al predicat, eliminable, que afirma
que \(m\) és una imatge de \(f\) amb aplicació exhaustiva associada
\(e\). Desplegat, aquest predicat quantifica el domini \(I\), l'aplicació
\(e\), i les condicions que \(m\) és un monomorfisme, que \(e\) és un
epimorfisme, que \(f=m\circ e\) i la minimalitat. Com és habitual, un cop
introduïdes les abreviacions les farem servir en les proves; d'ara
endavant ometrem aquest desplegament quan els tipus siguin inequívocs.
\end{notacio}

La definició demana la minimalitat entre totes les parts a través de les
quals passa \(f\). El resultat principal del capítol és que aquesta part
mínima existeix sempre i que es pot construir explícitament: és el
quocient de \(A\) pel parell nucli de \(f\).

\begin{teorema}[Existència i factorització]\label{teo:factoritzacio}
Tota aplicació \(f\colon A\to B\) admet una imatge. A més, existeix una
factorització
\[
A\overset{e}{\twoheadrightarrow}\operatorname{Im}(f)\overset{m}{\hookrightarrow}B,
\]
on \(e\) és un epimorfisme regular i \(m\) un monomorfisme, i aquesta
factorització és única llevat d'un únic isomorfisme compatible.
\end{teorema}

\begin{prova}
\emph{Construcció.} Prenem el parell nucli \(n=\operatorname{Nuc}(f)\),
que és una relació d'equivalència sobre \(A\)
(proposició~\ref{prop:kernel-equivalencia}), i el seu quocient
\(e\colon A\twoheadrightarrow A/n\)
(proposició~\ref{prop:quocient-existeix}), un epimorfisme regular
(proposició~\ref{prop:quocient-regular}). Per definició de parell nucli,
\(f\circ k_1=f\circ k_2\), on \(k_1,k_2\) són les projeccions de \(n\); és
a dir, \(f\) coiguala el parell \((k_1,k_2)\). Com que \(e\) és el
coequalitzador d'aquest parell, la seva propietat universal
(definició~\ref{def:coequalitzador}) proporciona una única aplicació
\(m\colon A/n\to B\) amb
\[
f=m\circ e.
\]
Anomenem \(I:=A/n\).

\emph{\(m\) és un monomorfisme.} Per l'observació~\ref{obs:injectiva-mono}
del capítol~\ref{cap:parts}, que usa la separació global, n'hi ha prou amb
veure que \(m\) és injectiva sobre els punts globals. Siguin
\(y,y'\colon\mathbf1\to I\) amb \(m\circ y=m\circ y'\). Com que \(e\) és un
epimorfisme, és exhaustiu sobre punts
(proposició~\ref{prop:epi-exhaustiu-punts}): existeixen
\(a,a'\colon\mathbf1\to A\) amb \(e\circ a=y\) i \(e\circ a'=y'\).
Aleshores
\[
f\circ a=m\circ e\circ a=m\circ y=m\circ y'=m\circ e\circ a'=f\circ a',
\]
de manera que el parell \((a,a')\) es factoritza pel parell nucli \(n\)
de \(f\): hi ha \(d\colon\mathbf1\to A\times_BA\) amb \(k_1\circ d=a\) i
\(k_2\circ d=a'\). Com que \(e\) coiguala \(k_1\) i \(k_2\),
\[
y=e\circ a=e\circ k_1\circ d=e\circ k_2\circ d=e\circ a'=y'.
\]
Per tant \(m\) és injectiva sobre punts i, doncs, un monomorfisme.

\emph{Minimalitat.} Suposem que \(f\) es factoritza també a través d'una
part \(m'\colon I'\hookrightarrow B\), diguem \(f=m'\circ g\) amb
\(g\colon A\to I'\). Aleshores \(m'\circ g\) coiguala el parell nucli de
\(f\): en efecte, \(m'\circ g\circ k_1=f\circ k_1=f\circ k_2=m'\circ
g\circ k_2\), i com que \(m'\) és un monomorfisme, \(g\circ k_1=g\circ
k_2\). Per la propietat universal del coequalitzador \(e\), \(g\) es
factoritza de manera única com \(g=t\circ e\) per a alguna
\(t\colon I\to I'\). Aleshores \(m'\circ t\circ e=m'\circ g=f=m\circ e\);
com que \(e\) és un epimorfisme, \(m'\circ t=m\). Això mostra que \(m\) es
factoritza a través de \(m'\), que és la minimalitat.

\emph{Unicitat.} Siguin \(m_1\colon I_1\hookrightarrow B\) i
\(m_2\colon I_2\hookrightarrow B\) dues imatges de \(f\), amb aplicacions
exhaustives \(e_1,e_2\). Per la minimalitat de cadascuna respecte de
l'altra, \(m_1\) es factoritza a través de \(m_2\) i recíprocament: hi ha
\(u\colon I_1\to I_2\) i \(v\colon I_2\to I_1\) amb \(m_2\circ u=m_1\) i
\(m_1\circ v=m_2\). Aleshores \(m_1\circ v\circ u=m_2\circ u=m_1\); com
que \(m_1\) és un monomorfisme, \(v\circ u=\id_{I_1}\), i simètricament
\(u\circ v=\id_{I_2}\). Per tant \(u\) és un isomorfisme, i és compatible
amb les factoritzacions perquè \(u\circ e_1=e_2\) se'n dedueix component
\(m_2\circ u\circ e_1=m_1\circ e_1=f=m_2\circ e_2\) i cancel·lant el
monomorfisme \(m_2\). La unicitat d'\(u\) és clara pel mateix argument de
cancel·lació.

\emph{Tota factorització en epimorfisme regular i monomorfisme és una
imatge.} Falta veure que la unicitat s'aplica a \emph{qualsevol}
factorització \(f=m''\circ e''\) amb \(e''\colon A\twoheadrightarrow I''\)
epimorfisme regular i \(m''\colon I''\hookrightarrow B\) monomorfisme, i no
només a les que són mínimes per hipòtesi. Primer, \(e''\) és el
coequalitzador del seu propi parell nucli: si \(e''\) és el coequalitzador
d'un parell \(f',g'\colon C\to A\), aleshores \(e''\circ f'=e''\circ g'\) i
el parell \((f',g')\) es factoritza pel parell nucli de \(e''\) (propietat
universal del pullback); per tant tota aplicació que coiguala el parell
nucli coiguala \((f',g')\) i es factoritza de manera única per \(e''\), i
\(e''\) mateix coiguala el seu parell nucli. Segon, com que \(m''\) és un
monomorfisme, \(e''\circ u=e''\circ v\) si i només si \(f\circ u=f\circ v\),
de manera que el pullback de \(f\) amb ella mateixa ho és també de \(e''\)
amb ella mateixa: \(e''\) és un coequalitzador de les projeccions
\(k_1,k_2\) del parell nucli de \(f\). L'argument de la minimalitat,
repetit amb \(e''\) en lloc de \(e\), mostra aleshores que \(m''\) es
factoritza per tota part per la qual es factoritza \(f\): \(m''\) és una
imatge de \(f\), i el paràgraf anterior dona l'únic isomorfisme compatible
amb \((m,e)\).
\end{prova}

\begin{observacio}[Sense elecció]\label{obs:factoritzacio-eleccio}
La construcció de \(e\) i \(m\), la minimalitat i la unicitat no depenen
de cap axioma més enllà dels coequalitzadors i l'efectivitat. El pas que
podria semblar necessitar l'elecció ---que \(m\) sigui un monomorfisme---
s'ha fet amb la separació global, reduint-lo als punts globals, i amb
l'exhaustivitat dels epimorfismes sobre punts, que és un teorema del
capítol~\ref{cap:valors}. Amb l'elecció l'argument és encara més curt:
una secció \(s\) de \(e\) dona directament, per a \(x,x'\colon X\to I\)
amb \(m\circ x=m\circ x'\), que \(s\circ x\) i \(s\circ x'\) estan
relacionats pel parell nucli i, doncs, \(x=e\circ s\circ x=e\circ s\circ
x'=x'\); però no cal. En les teories sense separació global (com les
dels conjunts variables del capítol~\ref{cap:models}) cap de les dues
vies no és disponible, i la prova passa per elements generalitzats i
recobriments epimòrfics; és el preu de no tenir prou punts.
\end{observacio}

\begin{corollari}[Tot epimorfisme és regular; sense elecció]\label{cor:epi-regular-sense-eleccio}
Tota aplicació \(f\colon A\to B\) que és un epimorfisme és un epimorfisme
regular: és el coequalitzador del seu parell nucli.
\end{corollari}

\begin{prova}
Factoritzem \(f=m\circ e\) pel teorema~\ref{teo:factoritzacio}. Com que
\(f\) és un epimorfisme i \(f=m\circ e\), també \(m\) ho és; i \(m\) és un
monomorfisme. Per tant \(m\) és un bimorfisme i, per la
proposició~\ref{prop:bimorfisme-isomorfisme} del capítol~\ref{cap:parts},
un isomorfisme. Així \(f=m\circ e\) és el compost d'un coequalitzador amb
un isomorfisme, i és un coequalitzador del mateix parell: el parell
nucli de \(f\), que coincideix amb el de \(e\) pel
lema~\ref{lema:nucli-invariant}.
\end{prova}

\begin{observacio}[Què afegeix l'elecció, revisat]\label{obs:epi-regular-revisat}
Al capítol~\ref{cap:relacions} vam obtenir que tot epimorfisme és regular
com a conseqüència de l'elecció (proposició~\ref{prop:epi-regular-eleccio}).
El corol·lari anterior mostra que l'elecció hi era supèrflua: la
factorització imatge i el teorema del bimorfisme basten. El que l'elecció
continua afegint, i que no se segueix d'aquí, és que tot epimorfisme
sigui \emph{escindit}. Tenim, doncs, tres nocions en ordre de força:
epimorfisme, que amb aquest corol·lari coincideix amb epimorfisme regular;
i epimorfisme escindit, estrictament més fort sense l'elecció i
equivalent als anteriors amb ella.
\end{observacio}

\begin{observacio}[Lectura clàssica]\label{obs:imatge-classica}
En el model clàssic, la imatge de \(f\colon A\to B\) és el subconjunt
\(\{\,f(a):a\in A\,\}\subseteq B\) dels valors efectivament assolits. La
factorització \(A\twoheadrightarrow\operatorname{Im}(f)\hookrightarrow B\)
és la descomposició habitual d'una funció en la sobrejecció sobre la seva
imatge seguida de la inclusió de la imatge dins del codomini. El parell
nucli identifica els elements amb la mateixa imatge, i el quocient és el
conjunt de valors assolits, cadascun comptat una sola vegada.
\end{observacio}

\begin{lema}[Estabilitat dels epimorfismes per canvi de base]\label{lema:estabilitat-epi-regular}
Si \(e\colon A\twoheadrightarrow C\) és un epimorfisme i \(g\colon D\to C\)
és una aplicació qualsevol, aleshores la projecció
\(e'\colon D\times_C A\to D\) del pullback de \(e\) al llarg de \(g\) és
també un epimorfisme (i, pel corol·lari~\ref{cor:epi-regular-sense-eleccio},
un epimorfisme regular). No s'usa l'elecció.
\end{lema}

\begin{prova}
Considerem el pullback
\[
\begin{aligned}
D\times_C A &\xrightarrow{\;\pi_A\;} A\\
{\scriptstyle e'}\big\downarrow\ \ &\qquad\big\downarrow{\scriptstyle e}\\
D &\xrightarrow{\;\;g\;\;} C.
\end{aligned}
\]
Per la proposició~\ref{prop:epi-exhaustiu-global} del
capítol~\ref{cap:relacions}, n'hi ha prou amb veure que \(e'\) és
exhaustiva sobre punts. Sigui \(d\colon\mathbf1\to D\). El punt
\(g\circ d\colon\mathbf1\to C\) té una preimatge per \(e\), que és
exhaustiu sobre punts: \(a\colon\mathbf1\to A\) amb \(e\circ a=g\circ d\).
El parell \((d,a)\) és compatible, i la propietat universal del pullback
dona \(w\colon\mathbf1\to D\times_CA\) amb \(e'\circ w=d\). Per tant tot
punt de \(D\) té preimatge per \(e'\).

Amb l'elecció hi ha una prova que no passa pels punts: una secció \(s\)
de \(e\) dona la secció \(\langle\id_D,s\circ g\rangle\) de \(e'\), de
manera que \(e'\) és fins i tot escindit.
\end{prova}

\begin{observacio}[Una teoria regular]\label{obs:regularitat}
La factorització del teorema~\ref{teo:factoritzacio} és el que en la
terminologia categòrica es demana a una categoria \emph{regular}: tota
aplicació es descompon com un epimorfisme regular seguit d'un
monomorfisme, i, pel lema~\ref{lema:estabilitat-epi-regular}, els
epimorfismes regulars són estables per canvi de base. Amb l'axioma
d'efectivitat del capítol~\ref{cap:relacions} ja teníem les relacions
d'equivalència efectives; ara veiem que la factorització imatge n'és una
conseqüència i que, un cop demostrada l'estabilitat, la nostra teoria
disposa de la infraestructura d'una categoria regular~\cite{borceux-vol2},
i, com que tot epimorfisme és regular, d'una en què els epimorfismes i
els epimorfismes regulars coincideixen. Tot això s'ha obtingut sense
l'axioma d'elecció, però no sense la separació global, que és el que ha
permès reduir les proves als punts globals; la
secció~\ref{sec:imatge-directa} en complementa la lectura amb l'adjunció
entre imatge directa i imatge inversa.
\end{observacio}

\section{Caracteritzacions per la factorització}\label{sec:caracteritzacions}

La factorització imatge unifica en un sol enunciat les nocions de
monomorfisme, epimorfisme i isomorfisme: cadascuna correspon a una
degeneració d'un dels dos factors.

\begin{proposicio}[Monomorfismes i parell nucli trivial]\label{prop:mono-nucli-trivial}
Una aplicació \(f\colon A\to B\) és un monomorfisme si i només si el seu
parell nucli és la diagonal, és a dir,
\(\operatorname{Nuc}(f)\equiv_{A\times A}\Delta_A\). En aquest cas
l'aplicació exhaustiva de la factorització és un isomorfisme
\(A\xrightarrow{\sim}\operatorname{Im}(f)\).
\end{proposicio}

\begin{prova}
Si \(f\) és un monomorfisme, aleshores \(f\circ k_1=f\circ k_2\) implica
\(k_1=k_2\), de manera que el parell nucli és la diagonal
(proposició~\ref{prop:diagonal-mono}). Recíprocament, si
\(\operatorname{Nuc}(f)\equiv_{A\times A}\Delta_A\) i \(x,x'\colon X\to A\)
compleixen \(f\circ x=f\circ x'\), aleshores \(\langle x,x'\rangle\) es
factoritza a través del parell nucli, per tant a través de la diagonal, i
això dona \(x=x'\) (proposició~\ref{prop:diagonal-mono}); per tant \(f\) és
un monomorfisme. En aquest cas el quocient de \(A\) per la diagonal és el
mateix \(A\), de manera que \(e\colon A\to\operatorname{Im}(f)\) és un
isomorfisme.
\end{prova}

\begin{proposicio}[Epimorfismes i imatge total]\label{prop:epi-imatge-total}
Una aplicació \(f\colon A\to B\) és un epimorfisme (equivalentment, un
epimorfisme regular) si i només si la seva imatge és tot el codomini, és
a dir, si el monomorfisme \(m\colon\operatorname{Im}(f)\hookrightarrow B\)
és un isomorfisme.
\end{proposicio}

\begin{prova}
Si \(m\) és un isomorfisme, \(f=m\circ e\) és el compost d'un epimorfisme
regular amb un isomorfisme, i per tant un epimorfisme regular.
Recíprocament, si \(f\) és un epimorfisme, \(m\) és un epimorfisme (perquè
\(f=m\circ e\)) i un monomorfisme, i per tant un isomorfisme pel teorema
del bimorfisme (proposició~\ref{prop:bimorfisme-isomorfisme}); és
l'argument del corol·lari~\ref{cor:epi-regular-sense-eleccio}.
\end{prova}

\begin{corollari}[Bimorfisme i isomorfisme]\label{cor:bimorfisme-iso}
Una aplicació és un isomorfisme si i només si és alhora un monomorfisme i
un epimorfisme. Equivalentment, si i només si tots dos factors de
la seva factorització imatge són isomorfismes.
\end{corollari}

\begin{prova}
Si \(f\) és un isomorfisme, és tant mono com epi. Recíprocament,
si \(f\) és un monomorfisme, l'aplicació exhaustiva \(e\) és un
isomorfisme (proposició~\ref{prop:mono-nucli-trivial}); si a més és un
epimorfisme, el monomorfisme \(m\) també ho és
(proposició~\ref{prop:epi-imatge-total}). Aleshores \(f=m\circ e\) és
composició de dos isomorfismes, i és un isomorfisme. El resultat ja es va
demostrar al capítol~\ref{cap:parts} amb el classificador de parts; el que
oferim aquí és la seva lectura a través de la factorització imatge: cada
factor recull una de les dues hipòtesis.
\end{prova}

\section{Imatge directa d'una part}\label{sec:imatge-directa}

La factorització imatge s'aplica en particular a la composició d'una part
amb una aplicació, i produeix la noció d'\emph{imatge directa}, dual de la
imatge inversa del capítol~\ref{cap:parts}.

\begin{definicio}[Imatge directa]\index{Imatge directa}\label{def:imatge-directa}
Siguin \(f\colon A\to B\) una aplicació i \(s\colon S\hookrightarrow A\)
una part de \(A\). La \emph{imatge directa} de \(s\) per \(f\) és la
imatge de la composició \(f\circ s\colon S\to B\):
\[
f_{*}(s):=\operatorname{Im}(f\circ s)\;\colon\; f_{*}S\hookrightarrow B.
\]
\end{definicio}

\begin{notacio}[Imatge directa]\label{not:imatge-directa}
Escriurem \(f_{*}(s)\) o \(f_{*}S\) per a la part de \(B\) obtinguda, i
reservarem \(f^{*}\) per a la imatge inversa (definició~\ref{def:imatge-inversa}).
El predicat eliminable corresponent és
\(\operatorname{Imatge}(m,e,f\circ s)\) de la notació~\ref{not:imatge}
aplicat a \(f\circ s\).
\end{notacio}

\begin{proposicio}[Adjunció imatge directa\,--\,imatge inversa]\label{prop:adjuncio-imatges}
Siguin \(f\colon A\to B\), una part \(s\colon S\hookrightarrow A\) i una
part \(t\colon T\hookrightarrow B\). Aleshores
\[
f_{*}S\subseteq_B T
\qquad\text{si i només si}\qquad
S\subseteq_A f^{*}T.
\]
\end{proposicio}

\begin{prova}
Suposem \(f_{*}S\subseteq_B T\). Sobre elements generalitzats, un
\(x\colon X\to A\) es factoritza a través de \(s\) precisament quan
pertany a \(S\); en aquest cas \(f\circ x\) es factoritza a través de
\(f\circ s\) i, per la definició d'imatge, a través de \(f_{*}S\), i per
hipòtesi a través de \(t\). Per tant \(f\circ x\) pertany a \(T\), és a
dir, \(x\) pertany a la imatge inversa \(f^{*}T\)
(proposició~\ref{prop:pertinenca-imatge-inversa}); això és \(S\subseteq_A
f^{*}T\). Recíprocament, si \(S\subseteq_A f^{*}T\), tot element
generalitzat de \(S\) té imatge per \(f\) dins de \(T\); com que
\(f_{*}S\) és la part mínima de \(B\) a través de la qual passa \(f\circ
s\), i \(t\) és una part a través de la qual passa \(f\circ s\), la
minimalitat dona \(f_{*}S\subseteq_B T\).
\end{prova}

\begin{observacio}[Lectura clàssica]\label{obs:imatge-directa-classica}
En conjunts, la imatge directa és \(f_{*}(S)=\{\,f(a):a\in S\,\}\) i la
inversa és \(f^{*}(T)=\{\,a:f(a)\in T\,\}\). La
proposició~\ref{prop:adjuncio-imatges} és l'equivalència elemental
\(f(S)\subseteq T\iff S\subseteq f^{-1}(T)\), que expressa que prendre
imatge directa i prendre imatge inversa són operacions adjuntes entre les
parts de \(A\) i les de \(B\). Estructuralment, aquesta adjunció no
necessita cap axioma nou: només la factorització imatge i la propietat
universal de la imatge inversa.
\end{observacio}

\begin{corollari}[La imatge total és la imatge de \(f\)]\label{cor:estabilitat-exhaustiva}
La imatge directa de la part total de \(A\) és la imatge de \(f\):
\(f_{*}A=\operatorname{Im}(f)\). En conseqüència, \(f\) és un epimorfisme
si i només si \(f_{*}A\equiv_B\id_B\).
\end{corollari}

\begin{prova}
La part total de \(A\) és \(\id_A\), i \(f\circ\id_A=f\), de manera que
\(f_{*}A=\operatorname{Im}(f)\). La segona afirmació és la
proposició~\ref{prop:epi-imatge-total}.
\end{prova}

\section{Composició de relacions}\label{sec:composicio-relacions}

Les eines d'aquest capítol permeten compondre relacions, i no només
aplicacions. Una relació de \(A\) cap a \(B\) i una relació de \(B\) cap a
\(C\) es componen prenent la imatge de la relació producte a través de la
projecció adequada. Recuperem així, estructuralment, la composició
clàssica de relacions.

\begin{definicio}[Composició de relacions]\index{Composició de relacions}\label{def:composicio-relacions}
Siguin \(r\colon R\hookrightarrow A\times B\) una relació de \(A\) cap a
\(B\) i \(r'\colon R'\hookrightarrow B\times C\) una relació de \(B\) cap
a \(C\). El seu \emph{producte fibrat sobre \(B\)} és el pullback de la
segona projecció de \(r\) amb la primera de \(r'\):
\[
R\times_B R'\hookrightarrow R\times R',
\]
que recull els parells de testimonis \((\rho,\rho')\) amb el segon
component de \(\rho\) igual al primer de \(\rho'\); és a dir, el pullback
de \(r_2\colon R\to B\) amb \(r'_1\colon R'\to B\), amb projeccions
\(\pi\colon R\times_B R'\to R\) i \(\pi'\colon R\times_B R'\to R'\) que
compleixen \(r_2\circ\pi=r'_1\circ\pi'\). La \emph{composició}
\(r'\circ r\) és la imatge de l'aplicació
\[
g:=\langle r_1\circ\pi,\;r'_2\circ\pi'\rangle\colon
R\times_B R'\longrightarrow A\times C,
\]
que a cada testimoni li assigna el primer component de \(\rho\) i el segon
de \(\rho'\).
\end{definicio}

Abans de llegir la composició per elements generalitzats necessitem un fet
que s'usa sovint quan es treballa amb recobriments: la pertinença a una
part es pot comprovar després d'un epimorfisme.

\begin{lema}[Pertinença després d'un recobriment]\label{lema:pertinenca-recobriment}
Siguin \(m\colon S\hookrightarrow B\) una part, \(y\colon X\to B\) un
element generalitzat i \(p\colon X'\twoheadrightarrow X\) un epimorfisme.
Si \(y\circ p\in_{X'}m\), aleshores \(y\in_Xm\).
\end{lema}

\begin{prova}
Sigui \(m'\colon P\hookrightarrow X\) la imatge inversa \(y^{-1}(m)\), amb
l'altra projecció \(q\colon P\to S\), de manera que \(y\circ m'=m\circ q\).
Per la proposició~\ref{prop:pertinenca-imatge-inversa}, la hipòtesi
\(y\circ p\in_{X'}m\) equival a \(p\in_{X'}m'\): existeix
\(w\colon X'\to P\) amb \(m'\circ w=p\). Com que \(p\) és un epimorfisme i
es factoritza per \(m'\), també \(m'\) és un epimorfisme; i és un
monomorfisme. Pel teorema del bimorfisme
(proposició~\ref{prop:bimorfisme-isomorfisme}), \(m'\) és un isomorfisme,
de manera que \(\id_X\in_Xm'\) i, de nou per la
proposició~\ref{prop:pertinenca-imatge-inversa}, \(y\in_Xm\).
\end{prova}

\begin{proposicio}[La composició és una relació]\label{prop:composicio-es-relacio}
Amb les notacions anteriors, \(r'\circ r\) és una relació de \(A\) cap a
\(C\). Sobre elements generalitzats, \(x\mathrel{(r'\circ r)_X}z\) es
compleix si i només si existeix un recobriment epimòrfic
\(p\colon X'\twoheadrightarrow X\) i un element \(y\colon X'\to B\) amb
\[
(x\circ p)\mathrel{r_{X'}}y
\qquad\text{i}\qquad
y\mathrel{r'_{X'}}(z\circ p).
\]
Amb l'axioma d'elecció, el recobriment \(p\) és escindit, de manera que el
testimoni es pot recollir ja en el mateix estadi \(X\).
\end{proposicio}

\begin{prova}
Que \(r'\circ r\) és una part de \(A\times C\) és immediat de la
definició d'imatge, que produeix sempre un monomorfisme cap al codomini
\(A\times C\) (teorema~\ref{teo:factoritzacio}). Provem la caracterització
sobre elements generalitzats construint explícitament el testimoni.

Factoritzem \(g=m\circ e\) amb \(e\colon R\times_B R'\twoheadrightarrow
\operatorname{Im}(g)\) epimorfisme regular i
\(m\colon\operatorname{Im}(g)\hookrightarrow A\times C\) monomorfisme, de
manera que \(r'\circ r=m\).

\emph{(\(\Leftarrow\)) De les components al testimoni de la composició.}
Suposem que hi ha un recobriment \(p\colon X'\twoheadrightarrow X\) i un
\(y\colon X'\to B\) amb \((x\circ p)\mathrel{r_{X'}}y\) i
\(y\mathrel{r'_{X'}}(z\circ p)\). Els dos testimonis són fletxes
\(\rho\colon X'\to R\) i \(\rho'\colon X'\to R'\) amb
\(r_1\circ\rho=x\circ p\), \(r_2\circ\rho=y\), \(r'_1\circ\rho'=y\) i
\(r'_2\circ\rho'=z\circ p\). Com que \(r_2\circ\rho=y=r'_1\circ\rho'\), la
propietat universal del pullback \(R\times_B R'\) (definició~\ref{def:pullback})
proporciona una única \(w\colon X'\to R\times_B R'\) amb \(\pi\circ w=\rho\)
i \(\pi'\circ w=\rho'\). Aleshores
\[
g\circ w=\langle r_1\circ\pi\circ w,\;r'_2\circ\pi'\circ w\rangle
=\langle r_1\circ\rho,\;r'_2\circ\rho'\rangle
=\langle x\circ p,\;z\circ p\rangle=\langle x,z\rangle\circ p,
\]
de manera que \(\langle x,z\rangle\circ p\) es factoritza a través de
\(g\), i per tant a través de \(m=r'\circ r\). Com que \(p\) és un
epimorfisme, el lema~\ref{lema:pertinenca-recobriment} dona que
\(\langle x,z\rangle\) pertany a la part \(m\), és a dir,
\(x\mathrel{(r'\circ r)_X}z\).

\emph{(\(\Rightarrow\)) Del testimoni de la composició a les components.}
Suposem \(x\mathrel{(r'\circ r)_X}z\), és a dir, que
\(\langle x,z\rangle\colon X\to A\times C\) es factoritza a través de
\(m\). Sense elecció, \(e\) no té per què ser exhaustiu sobre elements
generalitzats (proposició~\ref{prop:exhaustiva-generalitzada-escindit});
el que sí que val és que l'element
generalitzat \(\bar x\colon X\to\operatorname{Im}(g)\) que testimonia
aquesta factorització ---l'únic amb \(m\circ\bar x=\langle x,z\rangle\),
únic perquè \(m\) és un monomorfisme--- es realitza, sobre algun
recobriment epimòrfic \(p\colon X'\twoheadrightarrow X\), com
la imatge per \(e\) d'un element generalitzat
\(w\colon X'\to R\times_B R'\): basta prendre el pullback de \(e\) al
llarg de \(\bar x\), amb projeccions \(p\) ---un epimorfisme pel
lema~\ref{lema:estabilitat-epi-regular}--- i \(w\). És a dir,
\(e\circ w=\bar x\circ p\), d'on
\(g\circ w=m\circ e\circ w=m\circ\bar x\circ p=\langle x,z\rangle\circ
p\). Posem \(\rho:=\pi\circ w\),
\(\rho':=\pi'\circ w\) i \(y:=r_2\circ\rho\). La commutativitat del
pullback dona \(r'_1\circ\rho'=r'_1\circ\pi'\circ w=r_2\circ\pi\circ
w=r_2\circ\rho=y\). D'altra banda, la definició de \(g\) dona
\(r_1\circ\rho=x\circ p\) i \(r'_2\circ\rho'=z\circ p\). Per tant \(\rho\)
testimonia \((x\circ p)\mathrel{r_{X'}}y\) i \(\rho'\) testimonia
\(y\mathrel{r'_{X'}}(z\circ p)\), com volíem.

En tots dos casos, amb l'axioma d'elecció (capítol~\ref{cap:relacions}) tot
epimorfisme és escindit; prenent una secció de \(p\) es transporta el
testimoni al mateix estadi \(X\) i el recobriment es pot ometre.
\end{prova}

\begin{observacio}[Lectura clàssica]\label{obs:composicio-relacions-classica}
En conjunts, la composició de relacions és la relació habitual
\[
r'\circ r=\{\,(a,c):\exists b\,\bigl((a,b)\in r\text{ i }(b,c)\in
r'\bigr)\,\}.
\]
El quantificador existencial sobre \(b\) és exactament el que la imatge
introdueix: la projecció del producte fibrat recull tots els encadenaments
possibles, i la imatge els comprimeix en el conjunt de parells
\((a,c)\) efectivament connectats. En el llenguatge global, l'existència
del testimoni \(b\) per a cada parell pot requerir l'elecció; sobre
elements generalitzats, en canvi, el testimoni existeix sempre després
d'un canvi d'estadi.
\end{observacio}

\begin{definicio}[Relació funcional, total, i funcional i total]\index{Relació funcional, total, i funcional i total}\label{def:relacio-funcional}
Sigui \(r\colon R\hookrightarrow A\times B\) una relació de \(A\) cap a
\(B\), amb projeccions \(r_1\colon R\to A\) i \(r_2\colon R\to B\). Diem
que \(r\) és:
\begin{enumerate}
\item \emph{funcional} si la primera projecció \(r_1\colon R\to A\) és un
monomorfisme (com a molt un testimoni per a cada element de \(A\));
\item \emph{total} si la primera projecció \(r_1\colon R\to A\) és un
epimorfisme regular (com a mínim un testimoni per a cada element de
\(A\));
\item \emph{funcional i total} si és alhora funcional i total.
\end{enumerate}
\end{definicio}

Les dues condicions són independents i, juntes, equivalen a demanar que
\(r_1\) sigui un isomorfisme: un monomorfisme que és també un epimorfisme
regular és un isomorfisme (corol·lari~\ref{cor:bimorfisme-iso}). Convé
mantenir-les separades, perquè les relacions funcionals que no són totals
---les aplicacions \emph{parcials}--- apareixen de manera natural, i la
correspondència amb les aplicacions (proposició~\ref{prop:aplicacio-funcional})
usa precisament la totalitat com a condició addicional sobre la
funcionalitat.

Cadascuna de les dues meitats té la seva lectura sobre elements
generalitzats. La funcionalitat és la unicitat del testimoni; la
totalitat, la seva existència. Ho recollim per separat i després en
combinem les conseqüències.

\begin{proposicio}[Lectura de la funcionalitat i la totalitat]\label{prop:funcional-lectura}
Sigui \(r\colon R\hookrightarrow A\times B\) una relació. Aleshores:
\begin{enumerate}
\item \(r\) és funcional si i només si per a tot element generalitzat
\(x\colon X\to A\) hi ha \emph{com a molt} un \(y\colon X\to B\) amb
\(x\mathrel{r_X}y\);
\item \(r\) és total si i només si per a tot element generalitzat
\(x\colon X\to A\) existeix \emph{almenys} un \(y\colon X\to B\) amb
\(x\mathrel{r_X}y\), possiblement després d'un canvi d'estadi;
\item \(r\) és funcional i total si i només si per a tot \(x\colon X\to A\)
existeix un \emph{únic} \(y\colon X\to B\) amb \(x\mathrel{r_X}y\), ja en
el mateix estadi \(X\).
\end{enumerate}
\end{proposicio}

\begin{prova}
Recordem que \(x\mathrel{r_X}y\) (definició~\ref{def:relacio}) significa que la parella
\(\langle x,y\rangle\colon X\to A\times B\) es factoritza a través de la
part \(r\), és a dir, que existeix \(\rho\colon X\to R\) amb
\(r\circ\rho=\langle x,y\rangle\); aquest \(\rho\), quan existeix, és el
\emph{testimoni} de la relació. Component amb les projeccions del
producte, això equival a \(r_1\circ\rho=x\) i \(r_2\circ\rho=y\).

\emph{Funcionalitat.} Suposem \(r_1\) mono. Si \(x\mathrel{r_X}y\) i
\(x\mathrel{r_X}y'\), tenim testimonis \(\rho,\rho'\colon X\to R\) amb
\(r_1\circ\rho=x=r_1\circ\rho'\); com que \(r_1\) és un monomorfisme,
\(\rho=\rho'\), i per tant \(y=r_2\circ\rho=r_2\circ\rho'=y'\). Hi ha,
doncs, com a molt un \(y\). Recíprocament, suposem la unicitat del \(y\) i
siguin \(\rho,\rho'\colon X\to R\) amb \(r_1\circ\rho=r_1\circ\rho'=:x\).
Posant \(y:=r_2\circ\rho\) i \(y':=r_2\circ\rho'\), tots dos són
testimonis de \(x\), d'on \(y=y'\); llavors \(\rho\) i \(\rho'\) tenen les
dues components iguals, i com que \(r=\langle r_1,r_2\rangle\) és un
monomorfisme, \(\rho=\rho'\). Això prova que \(r_1\) és un monomorfisme.

\emph{Totalitat.} Suposem \(r_1\) epimorfisme, i sigui \(x\colon X\to A\).
Prenem el pullback de \(r_1\) al llarg de \(x\), amb projeccions
\(u\colon X'\to X\) i \(\rho\colon X'\to R\), de manera que
\(r_1\circ\rho=x\circ u\); pel lema~\ref{lema:estabilitat-epi-regular},
\(u\) és un epimorfisme. Posant \(y:=r_2\circ\rho\colon X'\to B\), el
testimoni \(\rho\) dona \((x\circ u)\mathrel{r_{X'}}y\): aquesta és
l'existència \emph{a menys del canvi d'estadi} \(u\). Recíprocament, si
per a \(x=\id_A\) existeix un testimoni després d'un recobriment
\(u\colon X'\twoheadrightarrow A\), tenim \(\rho\colon X'\to R\) amb
\(r_1\circ\rho=u\); com que \(u\) és un epimorfisme que es factoritza per
\(r_1\), també \(r_1\) és un epimorfisme, i regular pel
corol·lari~\ref{cor:epi-regular-sense-eleccio}.

\emph{Funcional i total.} És la conjunció de les dues anteriors: la
totalitat dona l'existència d'un testimoni i la funcionalitat en dona la
unicitat, de manera que el \(y\) existeix i és únic. Quan totes dues es
compleixen, \(r_1\) és un isomorfisme
(corol·lari~\ref{cor:bimorfisme-iso}) i el testimoni és
\(\rho=r_1^{-1}\circ x\), definit ja en el mateix estadi \(X\) sense
passar a cap recobriment; llavors \(y=r_2\circ r_1^{-1}\circ x\).
\end{prova}

\begin{proposicio}[Aplicacions com a relacions funcionals]\label{prop:aplicacio-funcional}
Assignar a cada aplicació \(f\colon A\to B\) el seu graf
\(\langle\id_A,f\rangle\colon A\hookrightarrow A\times B\) estableix una
correspondència biunívoca, llevat d'equivalència de parts, entre les
aplicacions de \(A\) cap a \(B\) i les relacions funcionals de \(A\) cap a
\(B\) que a més són totals ---és a dir, funcionals i totals en el sentit
de la definició~\ref{def:relacio-funcional}, amb primera projecció alhora
monomorfisme i epimorfisme regular---.
Sota aquesta correspondència, la composició de relacions de la
definició~\ref{def:composicio-relacions} coincideix amb la composició
d'aplicacions.
\end{proposicio}

\begin{prova}
El graf d'una aplicació és un monomorfisme perquè la primera projecció
del producte el reparametritza com \(\id_A\)
(proposició~\ref{prop:graf-part}); de fet la seva primera projecció
\emph{és} \(\id_A\), un isomorfisme, de manera que el graf és funcional i
total en el sentit de la definició~\ref{def:relacio-funcional}, amb
\(y=f\circ x\) l'únic element relacionat amb cada \(x\).
Recíprocament, una relació funcional i total té primera projecció
\(r_1\colon R\to A\) alhora monomorfisme i epimorfisme, i per tant un
isomorfisme (corol·lari~\ref{cor:bimorfisme-iso}); doncs
\(f:=r_2\circ r_1^{-1}\colon A\to B\) és una aplicació el graf de la qual
és equivalent a \(r\). Aquestes dues assignacions són mútuament inverses
llevat d'equivalència de parts. Que la composició de relacions es
correspon amb la d'aplicacions es comprova directament. Siguin
\(f\colon A\to B\) i \(g\colon B\to C\), amb grafs
\(r=\langle\id_A,f\rangle\) i \(r'=\langle\id_B,g\rangle\), de manera que
\(r_1=\id_A\), \(r_2=f\), \(r'_1=\id_B\) i \(r'_2=g\). El pullback de
\(r_2=f\) amb \(r'_1=\id_B\) és \(A\) mateix, amb projeccions
\(\pi=\id_A\) i \(\pi'=f\): si \(u\colon X\to A\) i \(v\colon X\to B\)
compleixen \(f\circ u=v\), l'única \(w\colon X\to A\) amb \(\id_A\circ
w=u\) i \(f\circ w=v\) és \(w=u\). Aleshores l'aplicació de la
definició~\ref{def:composicio-relacions} és
\(\langle r_1\circ\pi,\,r'_2\circ\pi'\rangle=\langle\id_A,g\circ f\rangle\),
que ja és un monomorfisme (proposició~\ref{prop:graf-part}); la seva
imatge és ella mateixa (proposició~\ref{prop:mono-nucli-trivial}), és a
dir, \(r'\circ r\) és el graf de \(g\circ f\).
\end{prova}

\begin{resumcap}
\apartat{Idees centrals}
\begin{enumerate}
  \item La imatge d'una aplicació és la part mínima del codomini per la qual es factoritza, i existeix sempre: és el quocient del domini pel parell nucli. La factorització \(f=m\circ e\), epimorfisme seguit de monomorfisme, és única llevat d'un únic isomorfisme.
  \item Que \(m\) sigui un monomorfisme es demostra sobre punts, amb la separació global i l'exhaustivitat dels epimorfismes sobre punts: no cal l'elecció.
  \item Conseqüència: tot epimorfisme és regular (coequalitzador del seu parell nucli), sense elecció. L'elecció només afegeix que sigui escindit. Els epimorfismes són estables per canvi de base, i la teoria és regular.
  \item Mono, epi i iso es llegeixen a la factorització: \(f\) és mono si \(e\) és iso (parell nucli diagonal), epi si \(m\) és iso (imatge total), iso si ho són tots dos.
  \item La imatge directa \(f_*S=\operatorname{Im}(f\circ s)\) és adjunta a la imatge inversa: \(f_*S\subseteq T\iff S\subseteq f^*T\).
  \item Les relacions es componen prenent la imatge del producte fibrat; sobre elements generalitzats, el testimoni intermedi existeix després d'un recobriment epimòrfic. Les aplicacions són exactament les relacions funcionals i totals, i la composició de relacions estén la d'aplicacions.
\end{enumerate}
\apartat{Errors típics} Pensar que la imatge és \enquote{el conjunt dels valors} com si fos un conjunt a part (és una part del codomini, un monomorfisme); creure que cal l'elecció per a la factorització imatge o per a \enquote{tot epi és regular}, quan basta la separació global; oblidar que la minimalitat és part de la definició d'imatge, i no només la factorització; confondre \(f_*\) (imatge directa, que necessita la factorització) amb \(f^*\) (imatge inversa, un pullback); suposar que el testimoni \(y\) de la composició \(r'\circ r\) existeix en el mateix estadi \(X\) sense elecció; i identificar \enquote{relació funcional} amb \enquote{aplicació}, quan falta la totalitat.
\apartat{Lectura recomanada} \cite[§§4.3--4.5 i cap.~5]{lawvere-rosebrugh} per a la factorització imatge, les imatges directes i la lectura de les relacions com a parts de productes; \cite[article~IV, sessions~22--23]{lawvere-schanuel} per a la imatge d'una aplicació amb exemples elementals; \cite[§§5.3--5.4 i §§9.6--9.7]{awodey} per a les categories regulars, les factoritzacions i la seva relació amb els límits; \cite[cap.~2]{borceux-vol2} per a la teoria general de les categories regulars i exactes, que és el marc del capítol; i \cite[§§7--10]{halmos} per a les relacions, les funcions, les famílies i les inverses en la teoria clàssica.
\end{resumcap}
